% Borrador.
\documentclass[11pt,a4paper]{amsart}

\usepackage[utf8]{inputenc}
\usepackage[T1]{fontenc}
\usepackage[spanish]{babel}
\usepackage{amsmath, amssymb, amsthm}
\usepackage{booktabs}
\usepackage{hyperref}
\hypersetup{hidelinks}
\usepackage{geometry}
\geometry{margin=2.5cm}

\newtheorem{teorema}{Teorema}
\newtheorem{corolario}[teorema]{Corolario}
\theoremstyle{remark}
\newtheorem{observacion}[teorema]{Observaci\'on}

\newcommand{\Z}{\mathbb{Z}}

\title[Un test de Miller verificado con la constante de Bach]{Un test de Miller verificado en Lean con la constante expl\'icita de Bach}
\author{Ernesto \'Alvarez}
\thanks{Este trabajo forma parte del proyecto AMate, de Ernesto \'Alvarez. Las demostraciones formales se desarrollaron con asistencia de modelos de lenguaje.}
\date{28 de noviembre de 2026 (borrador)}

\begin{document}

\begin{abstract}
Formalizamos en Lean~4, con Mathlib, el test de primalidad de Miller con una cota expl\'icita de tipo Bach: suponiendo GRH y varios enunciados anal\'iticos sobre funciones $L$ de Dirichlet, todo car\'acter no principal m\'odulo $n$ toma un valor distinto de $1$ en alg\'un $a\le 2(\log n)^2$ coprimo con $n$, si $\log n\ge 64$; con constantes $16$ y $4$ el umbral baja a $\log n\ge 8{,}3$ y $\log n\ge 23{,}5$. De ah\'i obtenemos un test de Miller ejecutable cuya correcci\'on est\'a demostrada para todo $n$, un programa de l\'inea de comandos que lo calcula, y los corolarios cl\'asicos de Bach: generadores peque\~nos de $(\Z/n\Z)^\times$, no-residuos de potencia $k$ peque\~nos y un test de Solovay--Strassen determinista. Las hip\'otesis anal\'iticas y GRH son teoremas de la biblioteca de Lean anunciada por Alp\"oge (2026); la composici\'on con ella deja todos estos resultados sin hip\'otesis. Todo depende solo de los tres axiomas est\'andar de Lean.
\end{abstract}

\maketitle

\section{Introducci\'on}

Miller~\cite{Miller} demostr\'o que, bajo la hip\'otesis de Riemann generalizada (GRH), basta probar las bases $a\le C(\log n)^2$ para decidir si $n$ es primo, y Bach~\cite{Bach} dio la constante $C=2$. Sin GRH, el test de Miller--Rabin~\cite{Rabin} es probabil\'istico. Lo que aportamos no es matem\'atica nueva, sino una formalizaci\'on completa en Lean~4~\cite{Lean4} con Mathlib~\cite{Mathlib}: la parte algebraica, la cota expl\'icita y un test que se ejecuta y cuya respuesta est\'a demostrada. La demostraci\'on de la constante es nuestra; sigue el esquema cl\'asico de la f\'ormula expl\'icita ponderada, pero no reproduce el art\'iculo de Bach l\'inea por l\'inea.

Las hip\'otesis que usamos son de dos tipos: GRH para caracteres primitivos y siete enunciados anal\'iticos sobre $L(s,\chi)$ (f\'ormula de Cauchy en rect\'angulos, localizaci\'on y finitud de los ceros, parte regular y ecuaci\'on funcional de $L'/L$, cotas de la funci\'on digamma y alturas buenas), m\'as la representaci\'on de la digamma como l\'imite. Todos son teoremas de la biblioteca de Lean anunciada por Alp\"oge~\cite{Alpoge}, verificada por m\'aquina~\cite{MIT} pero, hasta donde sabemos, sin revisi\'on humana publicada.

\section{La parte algebraica}

Sea $n\ge3$ impar. Un entero $a$ coprimo con $n$ es \emph{mentiroso fuerte} si $a^d\equiv1$ o $a^{2^rd}\equiv-1\pmod n$ para alg\'un $r<s$, donde $n-1=2^sd$ con $d$ impar. Si $n$ es primo, todo $a$ lo es; si $n$ es compuesto, los mentirosos fuertes est\'an contenidos en un subgrupo propio de $(\Z/n\Z)^\times$. Todo subgrupo propio est\'a contenido en el n\'ucleo de un car\'acter de Dirichlet no principal, por la dualidad de grupos abelianos finitos. Por tanto, si todo car\'acter no principal toma un valor distinto de $1$ en alg\'un $a\le B(n)$ coprimo con $n$, el test con las bases $a\le B(n)$ es correcto.

Lo mismo vale para el criterio de Euler, base del test de Solovay--Strassen~\cite{SS}: $a$ es \emph{mentiroso de Euler} si $a^{(n-1)/2}\equiv\left(\frac an\right)\pmod n$, con el s\'imbolo de Jacobi. Demostramos sin hip\'otesis que, para $n$ impar compuesto, los mentirosos de Euler forman un subgrupo propio (\texttt{primo\_de\_euler}): si $p^2\mid n$, una unidad de orden $p$ no cumple $u^{n-1}=1$; si no, $n=pm$ con $p\nmid m$, y por el teorema chino hay un $a$ que es no-residuo m\'odulo $p$ y $\equiv1\pmod m$, de modo que $\left(\frac an\right)=-1$ pero $a^{(n-1)/2}\equiv1\pmod m$.

\section{La cota de Bach con constante expl\'icita}

Sea $\chi\ne1$ un car\'acter m\'odulo $n$ y supongamos que $\chi(a)=1$ para todo $a\le x$ coprimo con $n$. Sean $\chi^*$ el car\'acter primitivo asociado, de conductor $q$, $L=\log n$ y $\ell=\log q$. Usamos el n\'ucleo $K(s)=x^s/((s+\alpha)(s+\beta))$ con $\alpha=2/5$, $\beta=9/20$, y el peso $w(m)=((m/x)^\alpha-(m/x)^\beta)/(\beta-\alpha)$ para $m<x$. La f\'ormula expl\'icita ponderada para $\chi^*$ y para $\zeta$, la cota de Hadamard de la suma sobre los ceros, cotas expl\'icitas de la digamma en el borde izquierdo, los residuos en $-\alpha$ y $-\beta$, y la contribuci\'on de las potencias de los primos que dividen a $n$ dan, para todo $x\ge1$ (\texttt{desigualdad\_bach}),
\begin{multline*}
\frac{100}{203}\,x\le\frac{\sqrt x}2\Bigl(1{,}082\,\ell+9{,}64\Bigr)+\frac{10\ell+672{,}9}{\sqrt x}\\
+\ell D(x)+24X(x)+51{,}3+4{,}01L+13{,}27(L-\ell),
\end{multline*}
con $D(x)=20(x^{-2/5}-x^{-9/20})$ y $X(x)=20(x^{-2/5}+x^{-9/20})$. Con $x=cL^2$ la desigualdad es falsa si $L$ es grande:

\begin{center}
\begin{tabular}{ccc}
\toprule
constante $c$ & umbral & teorema \\
\midrule
$2$ & $\log n\ge64$ & \texttt{cota\_bach\_explicita}\\
$4$ & $\log n\ge23{,}5$ ($n\ge2^{34}$) & \texttt{cota\_bach\_4}\\
$16$ & $\log n\ge8{,}3$ ($n\ge2^{12}$) & \texttt{cota\_bach\_16}\\
\bottomrule
\end{tabular}
\end{center}

\begin{teorema}
Bajo GRH y las hip\'otesis anal\'iticas, si $\log n\ge64$, todo car\'acter $\chi\ne1$ m\'odulo $n$ toma un valor distinto de $1$ en alg\'un $a\le2(\log n)^2$ coprimo con $n$; lo mismo con $a\le4(\log n)^2$ si $\log n\ge23{,}5$ y con $a\le16(\log n)^2$ si $\log n\ge8{,}3$.
\end{teorema}

\section{Consecuencias}

\begin{corolario}[Test certificado]
Sea $B(n)=n$ si $n<4096$; $8(\lfloor\log_2n\rfloor+1)^2$ si $4096\le n<2^{34}$; $2(\lfloor\log_2n\rfloor+1)^2$ si $2^{34}\le n<2^{93}$; y $(\lfloor\log_2n\rfloor+1)^2$ si $n\ge2^{93}$. Bajo las mismas hip\'otesis, el test de Miller con las bases $a\le B(n)$ decide la primalidad de todo $n$ (\texttt{esPrimoMiller\_\allowbreak certificado}), y tambi\'en el test de Euler con esas bases (\texttt{test\_euler\_\allowbreak certificado}).
\end{corolario}

\begin{corolario}[Generadores]
Si $\log n\ge64$, las unidades representadas por enteros $a\le2(\log n)^2$ generan $(\Z/n\Z)^\times$ (\texttt{generan\_explicito}); para todo $n\ge3$, lo hacen las representadas por $a\le B(n)$.
\end{corolario}

\begin{corolario}[No-residuos]
Sea $p$ primo y $k\ge2$ un divisor de $p-1$. Si $\log p\ge8{,}3$, hay un $a\le16(\log p)^2$ que no es potencia $k$-\'esima m\'odulo $p$, y uno con $a\le2(\log p)^2$ si $\log p\ge64$ (\texttt{noPotencia\_16}, \texttt{noPotencia\_explicito}). Con $k=2$ es la cota del menor no-residuo cuadr\'atico, en la l\'inea de Ankeny~\cite{Ankeny}.
\end{corolario}

Los tres corolarios salen de la misma observaci\'on: a todo subgrupo propio de $(\Z/n\Z)^\times$ le falta alg\'un elemento peque\~no.

\section{Composici\'on y ejecuci\'on}

Las hip\'otesis anal\'iticas son copias literales de enunciados de la biblioteca de Alp\"oge, y se instancian con sus teoremas por igualdad definicional (\texttt{HipotesisAlpoge.lean}, que no importa el m\'odulo principal y compila en pocos segundos). Con la GRH de esa biblioteca, los resultados anteriores quedan sin hip\'otesis. Compilamos la biblioteca desde las fuentes dos veces, en dos m\'aquinas virtuales distintas, y las dos veces sus $175\,187$ archivos \texttt{.olean} resultaron id\'enticos byte a byte a la lista del MIT~\cite{MIT}; sobre la segunda compilaci\'on compilaron las composiciones del test de Miller con constante $2$. Las composiciones posteriores (test para todo $n$ y no-residuos) compilaron localmente contra los binarios de la biblioteca. Todas dependen solo de \texttt{propext}, \texttt{Classical.choice} y \texttt{Quot.sound}, y el n\'ucleo independiente \texttt{nanoda} revisa sin errores el cierre de los teoremas principales.

El test es ejecutable. El programa \texttt{amate-primo} usa una copia de las funciones sin Mathlib, y \texttt{esPrimo\_eq} demuestra que calcula exactamente el test certificado. En un port\'atil decide $2^{127}-1$ ($16\,129$ bases) y otros diez n\'umeros en $1{,}6$ segundos en total, y $2^{521}-1$ ($271\,441$ bases) en $111$ segundos.

\section{Limitaciones}

\begin{itemize}
\item La versi\'on sin hip\'otesis depende de la biblioteca de Alp\"oge, verificada por m\'aquina pero sin revisi\'on humana publicada.
\item Las constantes no son \'optimas: Bach obtiene $2(\log n)^2$ para todo $n$, y nosotros solo para $\log n\ge64$.
\item La correcci\'on del programa se demuestra para las funciones de Lean; el c\'odigo nativo que genera el compilador no lo revisa el n\'ucleo.
\item El test de Euler no tiene versi\'on ejecutable, porque el s\'imbolo de Jacobi de Mathlib no est\'a pensado para calcular.
\end{itemize}

\begin{thebibliography}{9}
\bibitem{Alpoge} L.~Alp\"oge, nota de anuncio sin t\'itulo, \url{https://alpo.ge/remark.pdf}, 20 de noviembre de 2026. La nota no lleva firma; la atribuci\'on procede del dominio y de~\cite{MIT}.
\bibitem{Ankeny} N.~C.~Ankeny, \emph{The least quadratic non residue}, Ann.\ of Math.\ 55 (1952), 65--72.
\bibitem{Bach} E.~Bach, \emph{Explicit bounds for primality testing and related problems}, Math.\ Comp.\ 55 (1990), n\'um.~191, 355--380.
\bibitem{Lean4} L.~de~Moura, S.~Ullrich, \emph{The Lean~4 theorem prover and programming language}, CADE-28, LNCS 12699 (2021), 625--635.
\bibitem{Mathlib} The mathlib Community, \emph{The Lean mathematical library}, CPP 2020, 367--381.
\bibitem{Miller} G.~L.~Miller, \emph{Riemann's hypothesis and tests for primality}, J.\ Comput.\ System Sci.\ 13 (1976), 300--317.
\bibitem{MIT} A.~Sutherland, \emph{Verification files for the Lean proof of the Riemann Hypothesis}, \url{https://math.mit.edu/~drew/grh/}, consultado el 23 de noviembre de 2026.
\bibitem{Rabin} M.~O.~Rabin, \emph{Probabilistic algorithm for testing primality}, J.\ Number Theory 12 (1980), 128--138.
\bibitem{SS} R.~Solovay, V.~Strassen, \emph{A fast Monte-Carlo test for primality}, SIAM J.\ Comput.\ 6 (1977), 84--85.
\end{thebibliography}

\end{document}
